\documentclass{nsr}

\usepackage{amsmath}
\let\openbox\relax % 防御性处理，解决与 amsthm 的冲突
\usepackage{amsthm}
\usepackage{graphicx}
\usepackage{array}
\usepackage{makecell}
\usepackage{siunitx}
\usepackage{multirow}
\usepackage{booktabs}
\usepackage[table]{xcolor}
\usepackage{algorithm}
\usepackage{hyperref}
\sisetup{
  detect-weight=true,
  detect-inline-weight=math,
  input-symbols = (),
  table-number-alignment = center,
  table-format = 1.4
}

\definecolor{lightgray}{rgb}{0.88,0.88,0.88}

\makeatletter
\def\uns{\ifmmode\,\else$\,$\fi}%

\newtheorem{theorem}{Theorem}
\newtheorem{lemma}{Lemma}
\newtheorem{proposition}{Proposition}
\newtheorem{corollary}{Corollary}
\newtheorem{definition}{Definition}
\newtheorem{remark}{Remark}
\newtheorem{example}{Example}
\makeatother

\jvol{XX}
\jnum{X}
\jyear{Year}
\doi{10.1093/nsr/XXXX}
\received{XX XX Year}
\revised{XX XX Year}
\accepted{XX XX Year}

\markboth{One, Two, and Three}{One, Two, and Three}

\title{An Asynchronous Neuromorphic Architecture for Wearable EEG Emotion Recognition}

\author{Jingxin Liu$^{1,\dagger}$}
\author{Yiming Zhang$^{2,\dagger}$}
\author{Zikai Song$^{1}$}
\author{Zixuan Lai$^{2}$}
\author{Junyi Liu$^{1}$}
\author{Fuze Tian$^{3}$}
\author{Koo Hyoung Lee$^{4}$}
\author{Qunxi Dong$^{1,*}$}
\author{Lixian Zhu$^{1,*}$}
\author{Bin Hu$^{1,*}$}

\affil{$^1$Key Laboratory of Brain Health Intelligent Evaluation and Intervention (Ministry of Education), School of Medical Technology, Beijing Institute of Technology, Beijing 100081, China}
\affil{$^2$Key Laboratory of Brain Health Intelligent Evaluation and Intervention (Ministry of Education), Beijing Institute of Technology, Zhuhai 519088, China}
\affil{$^3$School of Information Science and Engineering, Lanzhou University, Lanzhou 730000, China}
\affil{$^4$NeuroSky, Inc., San Jose, CA 95112, USA}

\authornote{\textbf{Corresponding authors.} E-mail: dongqx@bit.edu.cn
zhulx@bit.edu.cn
bh@bit.edu.cn}

\authornote{Equally contributed to this work}

\begin{document}

\dhead{RESEARCH ARTICLE}
\subhead{INFORMATION SCIENCE}

\abstract[ABSTRACT]{
Affective electroencephalography (EEG) is event-driven and temporally sparse. However, prevailing synchronous, dense computing paradigms on general purpose hardware ignore this sparsity, incurring prohibitive energy costs that restrict long-term wearable monitoring. To address this, we propose a jointly optimized asynchronous neuromorphic architecture. Algorithmically, our dynamic adaptive spiking neural network model translates EEG dynamics into sparse events using a depthwise separable gating module for spatial efficiency, division-free time-to-first-spike encoding, and an adaptive temporal spiking neural network for stable learning. In hardware, these events drive a globally asynchronous locally synchronous accelerator, which aligns with the event-driven nature of EEG data, effectively minimizing global clock overheads and enabling computation only when required. Evaluated on the SEED dataset, the system achieves 96.98\% accuracy. Hardware evaluation shows dynamic power consumption of 40~mW and 2.66~$\mu$J per inference, providing an energy-efficient solution for continuous wearable deployment.
}

%\jelcode{Pa, J6, P16, E22}

\keywords{
affective electroencephalography, asynchronous neuromorphic architecture, spiking neural network, globally asynchronous locally synchronous, wearable monitoring}


\maketitle


\section{INTRODUCTION}\label{sec1}
Emotions are fundamental to human cognition, shaping perception, decision-making, and social interaction \cite{ref1}. Consequently, reliably and continuously decoding an individual's affective state in unconstrained environments has emerged as a critical frontier in brain-computer interfaces (BCIs) and personalized longitudinal healthcare \cite{ref2,ref3,ref4}. Among physiological modalities, electroencephalography (EEG) provides direct, high-temporal-resolution access to the central nervous system \cite{ref5,ref6,ref7} dynamics. Importantly, salient neural responses to affective stimuli are not continuous; rather, both microscopic neuronal spiking and macroscopic network reconfigurations in the brain are inherently event-driven and temporally sparse. This biological sparsity provides a principled foundation for highly efficient, on-demand cognitive tracking \cite{ref8}.

Recent progress in wearable EEG devices has markedly lowered the hardware barrier for long-term monitoring. Developments in lightweight sensors \cite{ref9}, hydrogel epidermal electrodes \cite{ref10}, and high-stability analog front-end circuits \cite{ref11} now enable prolonged EEG recording in daily life. However, signal acquisition alone is insufficient for truly wearable emotion recognition. In unconstrained environments, emotions are continuous, slowly varying, and strongly subject-dependent \cite{ref12}. These characteristics require the system to sustain reliable inference under strict constraints on power, intrusiveness, and continuous operation \cite{ref13}.

Convolutional neural networks (CNNs) have been widely adopted in emotion recognition and BCI tasks \cite{ref18}. For example, EEGNet \cite{ref56} typically processes continuously sampled EEG signals within fixed time windows to learn temporal dynamics and spatial patterns. However, its computational cost is tightly coupled to the input sampling rate; consequently, under high sampling rates or prolonged continuous monitoring, the overall overhead scales directly with the analysis duration. Long short-term memory (LSTM)-based models \cite{ref20,ref21} represent another common approach for EEG time series analysis, where sequence modeling is achieved through the recursive update of hidden states. Because these updates must adhere to temporal order, inference parallelization is severely constrained. Computation latency thus grows with sequence length, limiting both real-time performance and energy efficiency. Alternatively, graph convolutional network (GCN) methods, such as MSFR-GCN \cite{ref22}, utilize the spatial topology among electrodes for modeling. These approaches incorporate multi-scale spatial modeling and node interaction mechanisms, improving the representation of complex correlations across electrodes \cite{ref23,ref24}. While expressive, these models incur substantial overhead from graph aggregation and irregular access patterns, which stress the memory hierarchy.

Overall, the prevailing paradigm of synchronous acquisition and dense computation presents two key limitations. From an engineering perspective, its computational cost is tightly coupled to the input sampling rate, leading to substantial power and latency burdens during long-term continuous monitoring. From a neuroscience perspective, both microscopic neuronal spiking and macroscopic network reconfiguration in the brain are inherently distributed and asynchronous processes. Rigid global clock-triggering mechanisms are therefore ill-suited to capture the asynchronous temporal interactions characteristic of complex affective states.

Although model compression techniques such as pruning, quantization, and distillation reduce arithmetic costs to some extent \cite{ref25,ref26}, they fail to resolve the fundamental structural limitations of dense computation and global synchronization. This structural bottleneck becomes particularly acute for wearable EEG emotion recognition, which requires milliwatt power consumption and real-time responsiveness.
\begin{figure*}[t]
\centering
\includegraphics[width=0.95\textwidth]{Overview.eps}
\caption{\textbf{Jointly optimized neuromorphic architecture for asynchronous inference.}
EEG signals are transformed into spatial-spectral representations and processed by a lightweight dynamic adaptive spiking neural network model with depthwise separable convolution.
Dense activations are encoded into sparse spike times via division-free time-to-first-spike and executed on an address-event representation-based accelerator for energy-efficient inference.
GALS architecture is adopted at the system level to enable cross-domain asynchronous operation.}
\label{fig:system_overview}
\end{figure*}
Spiking neural networks (SNNs) offer a bio-inspired, event-driven alternative to conventional artificial neural networks (ANNs) \cite{ref27}. Because SNN neurons emit spikes only when their membrane potential crosses a specific threshold, they substantially reduce computation and memory traffic for sparse or slowly varying inputs \cite{ref28}. Several studies have demonstrated the feasibility of using SNNs for stress and emotion recognition on resource-constrained platforms \cite{ref29,ref30}. For example, Kumar et al. \cite{ref31} implemented an EEG decoding system using SNNs on neuromorphic processors to improve energy efficiency. Other work has shown that recurrent SNNs can effectively model the time evolution of EEG signals \cite{ref32}, and that simplified spiking dynamics enable depression recognition directly on microcontrollers \cite{ref33}. While these findings confirm that SNNs can exploit the inherent temporal sparsity of EEG data for efficient inference, realizing these energy benefits in wearable hardware remains challenging. Most current implementations run on conventional von Neumann processors, such as GPUs or general-purpose microcontrollers, which rely on global synchronous clocks. When event-driven spiking computation operates on synchronous hardware, the intended near-zero switching during inactive periods is undermined by always-on clocking and memory-centric data movement, largely offsetting the theoretical advantages of sparsity \cite{ref67,ref68}.

To resolve this paradigm mismatch, we propose a jointly optimized asynchronous neuromorphic architecture that connects temporally sparse brain dynamics with energy-efficient edge execution. At the core of our algorithmic framework is the Dynamic Adaptive Spiking Neural Network (DA-SNN), a lightweight spiking architecture featuring a depthwise separable gating module (DSGM) that decouples spatial and channel saliency to reduce spatiotemporal redundancy. To translate algorithmic sparsity into hardware-friendly events, we introduce a division-free time-to-first-spike (DF-TTFS) encoding coupled with an adaptive time-window mechanism. This design removes floating-point division from the inference-critical path and stabilizes temporal optimization. The resulting sparse event streams drive a custom globally asynchronous locally synchronous (GALS) accelerator \cite{ref69}. By replacing rigid global clocking with an asynchronous address-event handshake protocol, the system enables event-triggered, energy-proportional computation. This end-to-end co-design substantially reduces synchronization overhead, yielding an inference energy of 2.66~$\mu$J and 40~mW dynamic power while maintaining high accuracy across multiple benchmark datasets. This work provides a scalable cross-layer design template for continuous wearable cognitive interfaces and offers a paradigm readily transferable to future neuromorphic application-specific integrated circuit (ASIC) implementations.  


\begin{figure*}[t]
\hspace*{-2.8cm}
\begin{minipage}{\dimexpr\textwidth+2.8cm\relax}
\centering
\includegraphics[width=\textwidth]{DA-SNN.eps}
\caption{\textbf{The proposed DA-SNN model framework.} 
The model integrates DSGM for efficient spatiotemporal feature extraction, 
DF-TTFS encoder for sparse temporal coding, and an adaptive temporal spiking neural network for event-driven inference, enabling 
hardware-efficient EEG emotion recognition.}
\label{fig:DA-SNN}
\end{minipage}
\end{figure*}

\section{RESULTS}
In this section, we evaluate the proposed DA-SNN and its asynchronous hardware architecture. We first outline the neuromorphic framework for EEG emotion recognition, highlighting the synergy between the algorithm and event-driven hardware. We then analyze the learning dynamics of the spiking model and evaluate its classification performance across five standard EEG benchmarks. Finally, we assess the model's robustness against noise and low-precision quantization, followed by an analysis of the hardware's energy efficiency and physical field-programmable gate array (FPGA) implementation.


\subsection{Neuromorphic framework for EEG emotion recognition}

The proposed system-level framework for wearable EEG emotion recognition integrates data preprocessing, model architecture, and hardware deployment (Fig.~\ref{fig:system_overview}). The model architecture specifically encodes extracted emotion information into sparse event streams. To achieve this, we developed DA-SNN, an architecture integrating spatiotemporal feature extraction with event-driven spiking computation. As shown in Fig.~\ref{fig:DA-SNN}, the network processes data sequentially across three stages, beginning with DSGM for feature extraction, followed by TTFS encoding \cite{TTFS}, and concluding with an adaptive temporal spiking neural network (ATSNN).

The DSGM module combines depthwise separable convolution with parallel channel and spatial gating to capture discriminative features across all dimensions. The backbone applies depthwise separable convolution to the input feature map $X \in \mathbb{R}^{C\times H\times W}$:
\begin{equation}
F_{\mathrm{DS}} = f_{\mathrm{PW}}\!\bigl(f_{\mathrm{DW}}(X)\bigr)
\label{eq:dsgm_backbone}
\end{equation}
where $f_{\mathrm{DW}}(\cdot)$ and $f_{\mathrm{PW}}(\cdot)$ denote depthwise and pointwise convolutions, respectively, and $F_{\mathrm{DS}}$ is the backbone output. Concurrently, a dual-path gating mechanism refines these features. The channel gate aggregates global spatial information through average pooling to form a descriptor $s \in \mathbb{R}^{C \times 1 \times 1}$, producing the channel weights $G_{\mathrm{ch}} \in \mathbb{R}^{C \times 1 \times 1}$:
\begin{equation}
\left\{
\begin{aligned}
s_c &= \frac{1}{HW}\sum_{i=1}^{H}\sum_{j=1}^{W}X_{c,i,j},\quad c=1,\dots,C\\[0.5mm]
\sigma(x) &= \mathrm{clip}\!\left(\frac{x}{8}+0.5,\,0,\,1\right)\\[0.5mm]
G_{\mathrm{ch}} &= \sigma\!\left(f_{\mathrm{conv}}^{1\times1}(s_c)+b_{\mathrm{ch}}\right)
\end{aligned}
\right.
\label{eq:dsgm_ch_gate}
\end{equation}
where $b_{\mathrm{ch}}$ is the bias and $\sigma(\cdot)$ is the Hard-Sigmoid function. Similarly, the spatial gate compresses the channel dimension by computing the mean across all channels, yielding a 2D descriptor $M \in \mathbb{R}^{1 \times H \times W}$ and spatial weights $G_{\mathrm{sp}} \in \mathbb{R}^{1 \times H \times W}$:
\begin{equation}
\left\{
\begin{aligned}
M &= \frac{1}{C}\sum_{c=1}^{C}X_c\\[0.5mm]
G_{\mathrm{sp}} &= \sigma\!\left(f_{\mathrm{conv}}^{3\times3}(M)+b_{\mathrm{sp}}\right)
\end{aligned}
\right.
\label{eq:dsgm_sp_gate}
\end{equation}
where $b_{\mathrm{sp}}$ is the bias. The backbone features are then modulated by these learned weights through element-wise multiplication, followed by batch normalization and ReLU activation:
\begin{equation}
F_{\mathrm{out}} = f_{\mathrm{ReLU}}\!\Bigl(f_{\mathrm{BN}}\bigl(F_{\mathrm{DS}}\odot G_{\mathrm{ch}}\odot G_{\mathrm{sp}}\bigr)\Bigr)
\label{eq:dsgm_out}
\end{equation}
where $\odot$ denotes element-wise multiplication.

\begin{table*}[t]
\centering
\hspace*{-2.8cm}
\begin{minipage}{1.2\textwidth}
\caption{Performance Comparison with Classical Machine Learning and Lightweight Deep Learning Models}
\label{tab:overall_15models}
\setlength{\tabcolsep}{2.5pt}
\renewcommand{\arraystretch}{1.5}
\scriptsize
\resizebox{\textwidth}{!}{%
\begin{tabular}{l
S[table-format=2.2] S[table-format=2.2] S[table-format=2.2]
>{\columncolor{lightgray}}S[table-format=2.2] >{\columncolor{lightgray}}S[table-format=2.2] >{\columncolor{lightgray}}S[table-format=2.2]
S[table-format=2.2] S[table-format=2.2] S[table-format=2.2]
>{\columncolor{lightgray}}S[table-format=2.2] >{\columncolor{lightgray}}S[table-format=2.2] >{\columncolor{lightgray}}S[table-format=2.2]
S[table-format=2.2] S[table-format=2.2] S[table-format=2.2]}
\toprule
\multirow{2}{*}{Method}
& \multicolumn{3}{c}{SEED}
& \multicolumn{3}{>{\columncolor{lightgray}}c}{SEED-IV}
& \multicolumn{3}{c}{SEED-V}
& \multicolumn{3}{>{\columncolor{lightgray}}c}{DEAP}
& \multicolumn{3}{c}{DREAMER} \\
\cmidrule(lr){2-4}\cmidrule(lr){5-7}\cmidrule(lr){8-10}\cmidrule(lr){11-13}\cmidrule(lr){14-16}
& {Acc(\%)} & {F1(\%)} & {Spe(\%)}
& {Acc(\%)} & {F1(\%)} & {Spe(\%)}
& {Acc(\%)} & {F1(\%)} & {Spe(\%)}
& {Acc(\%)} & {F1(\%)} & {Spe(\%)}
& {Acc(\%)} & {F1(\%)} & {Spe(\%)} \\
\midrule
RF\cite{ref52} & 91.91 & 91.65 & 93.30 & 89.85 & 87.86 & 92.30 & 89.23 & 87.82 & 91.98 & 79.43 & 76.31 & 92.58 & 86.14 & 81.34 & 92.22 \\
k-NN\cite{ref53} & 90.82 & 89.90 & 92.12 & 89.05 & 87.71 & 92.26 & 87.91 & 87.24 & 91.27 & 82.75 & 79.86 & 90.20 & 88.88 & 88.02 & 91.89 \\
XGBoost\cite{ref54} & 88.11 & 86.60 & 90.72 & 86.27 & 85.37 & 91.35 & 88.41 & 85.20 & 92.30 & 75.26 & 75.47 & 91.21 & 85.56 & 85.46 & 89.95 \\
SqueezeNet\cite{ref57} & 94.04 & 93.31 & 96.58 & 92.35 & 92.63 & 97.44 & 90.37 & 91.29 & 97.68 & 67.74 & 73.00 & 85.38 & 91.08 & 90.17 & 96.64 \\
MobileNetV3-Small\cite{ref58} & 93.40 & 90.14 & 94.23 & 92.61 & 93.52 & 97.51 & 91.91 & 91.66 & 94.66 & 84.11 & 82.92 & 94.53 & 89.24 & 89.21 & 96.29 \\
MobileNetV3-Large\cite{ref58} & 90.85 & 91.76 & 92.66 & 90.34 & 88.88 & 94.87 & 89.73 & 88.36 & 88.94 & 86.29 & 82.59 & 92.17 & 91.63 & 91.37 & 97.62 \\
ShuffleNetV2\cite{ref59} & 91.01 & 92.43 & 95.51 & 93.71 & 92.61 & 97.91 & 90.22 & 91.22 & 95.47 & 77.21 & 77.17 & 92.02 & 92.85 & 92.88 & 97.49 \\
EfficientNet-B0\cite{ref60} & 95.35 & 94.60 & 97.67 & 92.16 & 92.34 & 97.34 & 91.65 & 93.34 & 97.58 & 83.92 & 83.90 & 94.43 & 88.35 & 87.54 & 94.04 \\
EfficientNet-B3\cite{ref60} & 96.03 & 95.58 & 98.02 & 92.96 & 93.58 & 97.60 & 91.19 & 91.35 & 96.96 & 70.14 & 72.71 & 89.35 & 85.54 & 86.09 & 92.56 \\
Deep ConvNet\cite{ref55} & 95.88 & 94.91 & 97.93 & 92.48 & 92.06 & 97.46 & 90.51 & 89.58 & 97.63 & 82.19 & 84.10 & 89.99 & 89.81 & 90.53 & 94.95 \\
Shallow ConvNet\cite{ref55} & 92.27 & 94.18 & 96.15 & 89.49 & 90.32 & 96.12 & 85.78 & 87.69 & 96.20 & 80.58 & 79.85 & 89.79 & 83.61 & 82.80 & 86.28 \\
EEGNet\cite{ref56} & 91.34 & 91.89 & 95.67 & 88.03 & 83.46 & 90.99 & 86.08 & 85.65 & 93.89 & 81.75 & 80.18 & 91.02 & 86.45 & 85.97 & 88.19 \\
EESCN\cite{ref61} & 96.08 & 94.85 & 98.04 & 92.60 & 92.54 & 97.15 & 89.01 & 90.15 & 95.91 & 87.65 & 86.55 & 94.82 & 92.32 & 91.84 & 97.27 \\
DH-SNN\cite{ref62} & 96.54 & 95.75 & 98.40 & 93.08 & 91.44 & 98.27 & 91.51 & 94.42 & 97.16 & 75.75 & 78.24 & 88.85 & 89.89 & 90.78 & 96.59 \\
\textbf{DA-SNN} & \textbf{96.98} & \textbf{96.91} & \textbf{98.49} & \textbf{93.87} & \textbf{94.10} & \textbf{98.35} & \textbf{92.10} & \textbf{93.31} & \textbf{97.72} & \textbf{88.09} & \textbf{88.04} & \textbf{95.26} & \textbf{94.31} & \textbf{94.88} & \textbf{98.21} \\
\bottomrule
\end{tabular}
}
\end{minipage}
\end{table*}

Conventional neuron models like Leaky Integrate-and-Fire (LIF) \cite{ref30} often present computational bottlenecks. We replace them with the B1 spiking neuron model \cite{ref46}, which uses a cascaded time-window mechanism to segment neuronal dynamics and simplify membrane potential evolution into a piecewise linear form. For layer $n$, dynamics are defined over a time window $[T_{\min}^{(n)}, T_{\max}^{(n)}]$, where $T_{\min}^{(n)} = T_{\max}^{(n-1)}$. The membrane potential $V_i^{(n)}$ of neuron $i$ follows a piecewise linear differential equation:
\begin{equation}
\label{eq:b1_dynamics}
\epsilon \frac{dV_i^{(n)}}{dt} =
\begin{cases}
\displaystyle \sum\nolimits_j W_{ij}^{(n)} H\!\left(t-T_j^{(n-1)}\right), & t < T_{\min}^{(n)}\\[0.4em]
1, & \hspace{-3.5em}T_{\min}^{(n)} \le t \le T_{\max}^{(n)}.
\end{cases}
\end{equation}
where $\epsilon$ is a scaling constant, and $H(\cdot)$ represents the Heaviside step function. Integrating Eq.~(\ref{eq:b1_dynamics}) provides a closed-form solution for the theoretical output spike time $\tilde{T}_i^{(n)}$:
\begin{equation}
\label{eq:b1_combined}
\left\{
\begin{aligned}
\tilde{T}_i^{(n)} &= T_{\min}^{(n)} + \epsilon \vartheta_i^{(n)} + \sum\nolimits_j W_{ij}^{(n)}\!\left(T_j^{(n-1)} - T_{\min}^{(n)}\right) \\
T_i^{(n)} &= \min\!\left(\tilde{T}_i^{(n)},\, T_{\max}^{(n)}\right)
\end{aligned}
\right.
\end{equation}
\noindent If the required integration time $\tilde{T}_{i}^{(n)}$ exceeds the upper bound, the output $T_i^{(n)}$ is clipped to $T_{\max}^{(n)}$, keeping the neuron silent during inference.

Continuous DSGM features are converted into sparse spike times for the downstream ATSNN using TTFS coding, which maps stronger activations to earlier spikes. While standard TTFS normalizes input activation $V_{\mathrm{in}}$ via division:
\begin{equation}
V_{\mathrm{norm}} = \frac{V_{\mathrm{in}} - V_{\min}}{V_{\max} - V_{\min}}
\label{eq:standard_ttfs}
\end{equation}
where $V_{\min}$ and $V_{\max}$ denote the minimum and maximum bounds, we implement a DF-TTFS encoder that adapts to the dynamic range. The complete encoding process is:  
\begin{equation}
\left\{
\begin{aligned}
S_p &= 2^{\lceil \log_2(V_{\max}-V_{\min}) \rceil} \\
V_{\mathrm{norm}} &= \mathrm{clip}\left(\frac{V_{\mathrm{in}}-V_{\min}}{S_p}, 0, 1\right) \\
T_{\mathrm{spike}} &= T_{\max} - \mathrm{round}\left(V_{\mathrm{norm}} (T_{\max}-T_{\min})\right)
\end{aligned}
\right.
\label{eq:ttfs_div_free}
\end{equation}
where $S_p$ is the scaling factor, $\mathrm{clip}(\cdot)$ clamps values to $[0,1]$, and $T_{\mathrm{spike}}$ is the discrete spike time.

Static time windows in SNNs using TTFS encoding can cause gradient instability. For a firing neuron, the gradient with respect to synaptic weights relies on presynaptic time offsets:
\begin{equation}
\frac{\partial L}{\partial W_{ij}^{(n)}} = \frac{\partial L}{\partial T_i^{(n)}} \cdot \left(T_j^{(n-1)} - T_{\min}^{(n)}\right)
\label{eq:grad_w}
\end{equation}
To resolve this, the ATSNN incorporates an adaptive temporal window adjustment mechanism. By monitoring the earliest valid spike time $T_{e}^{(n)}$ in layer $n$, the upper bound of the time window is updated:
\begin{equation}
T_{\max}^{(n)\prime} = T_{\max}^{(n)} + \lambda\!\left(T_{e}^{(n)} - \frac{T_{\max}^{(n)} + T_{\min}^{(n)}}{2}\right)
\label{eq:atsnn_update}
\end{equation}
where $\lambda$ controls the adaptation rate. This updated bound $T_{\max}^{(n)\prime}$ propagates to the subsequent layer by setting $T_{\min}^{(n+1)} = T_{\max}^{(n)\prime}$, making the time window dynamically adjustable (detailed in Supplementary Note~S2.3).

The output layer acts as a nonspiking integrator, accumulating spike contributions from the ATSNN within its time window. The logit for class $i$ is calculated as:
\begin{equation}
L_i = \sum\nolimits_j W_{ij}^{(L)}\!\left(T_{\min}^{(L)} - T_j^{(L-1)}\right) + \epsilon \vartheta_i^{(L)}
\label{eq:out_logit}
\end{equation}
The network predicts the class corresponding to the neuron with the highest logit value.


\subsection{Emotion Recognition Performance Across EEG Benchmarks}

We evaluate the proposed DA-SNN framework on five standard EEG emotion recognition datasets (dataset descriptions and the standardized preprocessing pipeline are detailed in Supplementary Methods~S1.1 and Methods~S1.2).

DA-SNN outperforms both classical machine learning methods and representative lightweight deep architectures across all five benchmarks (Table~\ref{tab:overall_15models}). On the SEED dataset, the model reaches 96.98\% accuracy, a 96.91\% F1 score, and 98.49\% specificity. This performance extends to tasks with finer emotional granularity, maintaining accuracies of 93.87\% and 92.10\% on the SEED-IV and SEED-V datasets, respectively. Whereas standard lightweight convolutional networks like SqueezeNet and EfficientNet exhibit variability across these finer-grained states, DA-SNN remains stable.

The model also demonstrates robustness under the variable, subject-dependent conditions of the DEAP and DREAMER datasets, achieving accuracies of 88.09\% and 94.31\%, respectively. Notably, DA-SNN achieves the highest specificity across all evaluated datasets, reflecting reliable negative sample identification.

We validate the specific contributions of individual architectural components, including the gating mechanisms and the adaptive temporal window, through systematic ablation studies (detailed in Supplementary Experiments~S3.1 and Table~S2).

\begin{figure*}[t]
\centering
\includegraphics[width=0.95\textwidth]{robustness_quantization_comparison.eps}
\caption{
\textbf{Noise robustness, quantization sensitivity, and model comparison.}
\textbf{(a),(b)} Results after 10 epochs of fine-tuning on the SEED dataset.
\textbf{(a)} Accuracy as a function of the standard deviation (SD) of injected random Gaussian noise.
\textbf{(b)} Quantization sensitivity when reducing weight precision to a given number of bits.
\textbf{(c)} Performance comparison of DA-SNN with representative lightweight models. FLOPs and parameter counts are normalized.
}
\label{fig:robust_quant_comparison}
\end{figure*}

\begin{table}[b]
\caption{Efficiency Analysis of Comparison Classifiers}
\label{tab:efficiency}

\setlength{\tabcolsep}{3pt} % 减少了列数，可以适当放宽列间距使表格更舒展
\renewcommand{\arraystretch}{1.08}
\scriptsize % 保持基础字号为 scriptsize

\sisetup{
  table-number-alignment=center,
  parse-numbers=true,
  detect-weight=true % 允许 S 列直接识别底部加粗的字体
}
% 将宽度限制为 0.96\columnwidth，右侧留出安全空隙，防止与右边分栏文本重叠
\resizebox{\columnwidth}{!}{%
\begin{tabular}{@{}l
S[table-format=3.2, table-space-text-post={~M}] % 为可能出现的 M 预留对齐空间
S[table-format=3.2, table-space-text-post={~M}]
S[table-format=3.2, table-space-text-post={~M}]
S[table-format=3.2, table-space-text-post={~MB}]
S[table-format=3.2, table-space-text-post={~MB}]
S[table-format=3.2]@{}}
\toprule
\textbf{Model}
& {\textbf{Params (K)}}
& {\textbf{FLOPs (K)}}
& {\textbf{\makecell{Mult-\\Adds (K)}}} 
& {\textbf{\makecell{Param\\Size (KB)}}} 
& {\textbf{RAM (KB)}} 
& {\textbf{Power (mW)}} \\
\midrule

SqueezeNet        &   4.12 & 259.97 & 129.98 &  16.09 &  50.43 & 368.33 \\
MobileNetV3-Small &  26.43 & 993.12 & 496.56 & 103.23 & 392.44 & 458.33 \\
MobileNetV3-Large & 543.90 &   9.69{~M} &   4.85{~M} &   2.07{~MB} &   3.53{~MB} & 546.35 \\
ShuffleNetV2      & 244.22 &   3.49{~M} &   1.74{~M} & 953.98 &   1.38{~MB} & 426.85 \\
EfficientNet-B0   &   1.07{~M} &  11.99{~M} &   6.00{~M} &   4.10{~MB} &   5.28{~MB} & 437.84 \\
EfficientNet-B3   &   1.38{~M} &  15.32{~M} &   7.66{~M} &   5.27{~MB} &   6.93{~MB} & 524.62 \\
Deep ConvNet      & 239.48 &   1.87{~M} & 936.20 & 935.46 & 995.98 & 435.53 \\
Shallow ConvNet   &  14.28 & 350.08 & 175.04 &  55.79 &  85.07 & 258.37 \\
EEGNet            &   2.72 &  87.55 &  43.78 &  10.64 &  41.04 & 334.86 \\
EESCN             &  11.08{~M} & 308.59{~M} & 154.29{~M} &  42.27{~MB} & 225.12 & 776.79 \\
DH-SNN            & 152.58 & 476.00 & 235.96 & 587.01 & 260.89 & 236.49 \\

\midrule
\bfseries DA-SNN (FP32)
 & \bfseries 12.76 
 & \bfseries 32.47 
 & \bfseries 16.23 
 & \bfseries 49.83 
 & \bfseries 51.27 
 & \bfseries 236.68 \\

\bfseries DA-SNN (INT8)
 & \bfseries 12.76 
 & {--} 
 & \bfseries 16.23 
 & \bfseries 12.50 
 & \bfseries 13.68  
 & \bfseries 184.14 \\
\bottomrule
\end{tabular}%
} 
\end{table}


\begin{figure*}[t]
\centering
\includegraphics[width=\textwidth]{hardoverview.eps}
\caption{\textbf{Hardware architecture and main components of the proposed accelerator.}
(a) System level organization integrating a streaming INT8 tensor compute engine, a DF-TTFS encoder, and event-driven SNN arrays within a GALS control framework.
(b) Line Buffer based on shift registers supporting continuous sliding window reconstruction for pipelined convolution.
(c) Hardware realization of the DF-TTFS encoder with bit shift normalization and event detection for address-event representation (AER) spike output.
(d) Microarchitecture of the event-driven processing element (PE) illustrating the time-modulated membrane potential integration and local asynchronous handshaking.}
\label{fig:hardoverview}
\end{figure*}

\subsection{Robustness Under Noise and Quantization}

Continuous EEG acquisition introduces motion artifacts, electrode impedance fluctuations, and environmental interference. To evaluate deployment reliability, we injected zero-mean Gaussian noise with varying standard deviations into the input feature maps after ten epochs of fine-tuning on the SEED dataset. The model demonstrates stable performance under mild perturbations (Fig.~\ref{fig:robust_quant_comparison}a). Increasing the noise standard deviation from 0.001 to 0.02 yields a minor accuracy decrease from 96.94\% to 96.09\%. Accuracy remains above 93\% even when the noise level reaches 0.05. Noticeable degradation only appears when the standard deviation exceeds 0.075, substantially reducing the signal-to-noise ratio. This noise tolerance arises from the event-driven TTFS encoding mechanism, which relies on relative activation magnitude rather than absolute amplitude. Because small additive perturbations cause only slight shifts in spike timing, the adaptive temporal window effectively regulates spike distributions to maintain stability under distributional shifts.  

We also examined sensitivity to weight quantization by reducing numerical precision from 32-bit floating-point to low-bit fixed-point representations. Decreasing precision to 8-bit introduces a negligible accuracy loss of 0.04\% (Fig.~\ref{fig:robust_quant_comparison}b), confirming that 8-bit integer deployment preserves discriminative performance. Accuracy remains above 95\% at 6-bit precision, indicating high tolerance to parameter discretization. Degradation only emerges at 4-bit precision, where limited representational granularity constrains synaptic resolution. Together, these results demonstrate that the model maintains stable performance under the signal interference and hardware constraints typical of practical deployment environments.

\subsection{Computational Efficiency and Energy Consumption}

We analyze the computational efficiency of the proposed architecture in terms of model scale, inference cost, and memory footprint, measuring practical power consumption directly on a Raspberry Pi 5 platform. The network contains only 12.76~K parameters (Table~\ref{tab:efficiency}), requiring significantly fewer resources than standard lightweight models such as MobileNetV3-Large and EfficientNet-B0 while preserving recognition performance.

Because the network executes event-driven inference using TTFS encoding, neurons fire at most once per cycle. This yields an empirical average firing rate of 35.26\%, which reduces the effective inference cost from a theoretical maximum of 32.47~K floating-point and 16.23~K multiply-accumulate operations down to 16.41~K and 8.21~K, respectively. The sparsity introduced by this temporal encoding minimizes runtime operations by restricting computation strictly to valid spike events. For the 32-bit floating-point model, parameter storage and runtime memory require 49.83~KB and 51.27~KB, respectively. Applying 8-bit integer quantization compresses these requirements to 12.50~KB and 13.68~KB. This quantized configuration meets stringent embedded constraints and achieves a measured energy consumption of just 261.21~$\mu$J per inference. Ultimately, the network balances recognition accuracy with low computational demand (Fig.~\ref{fig:robust_quant_comparison}c), satisfying the strict power and memory constraints of wearable deployment.



\subsection{Hardware Implementation}

To validate the practical deployability of our framework, we implemented a dedicated, energy-efficient heterogeneous hardware accelerator for edge EEG inference. To operate within the strict power and area limits of edge devices, this architecture uses an INT8 global quantization strategy alongside an asynchronous event-driven paradigm (Supplementary Note~S2.4). These design choices enable real-time inference while maintaining high accuracy under tight resource constraints. The system integrates a global low-bitwidth quantization datapath, a reconfigurable tensor-streaming compute subsystem, and a cross-domain asynchronous energy-saving mechanism.  

The hybrid streaming and event processing pipeline establishes an end-to-end data flow (Fig.~\ref{fig:hardoverview}a). Multi-channel data first stream into the on-chip input buffer via an SPI interface. The reconfigurable tensor compute engine (TCE) operates in the synchronous numeric domain, configuring its datapath in convolution mode to extract low-level features. To sustain continuous convolution throughput, the engine employs a shift-register-based line buffer that reconstructs sliding windows from the input stream for pipelined processing (Fig.~\ref{fig:hardoverview}b). After feature fusion through a cascaded systolic scheme, the fixed-point data stream enters a division-free TTFS encoder. This encoder converts dense representations into sparse temporal spike sequences, encapsulating them into AER packets containing addresses and timestamps \cite{AER}. These packets cross the clock-domain boundary into the asynchronous spiking domain through an asynchronous handshake protocol \cite{ref66}, where they drive the SNN compute units. Upon event arrival, the SNN engine activates only the required compute arrays to perform parallel weight retrieval and membrane potential accumulation. After membrane potential updates complete in the final neuron layer, the decision unit outputs the classification result to conclude the inference.

We introduce several hardware design optimizations to improve efficiency. The accelerator integrates a streaming reconfigurable TCE that uses time multiplexing on shared compute resources to support multiple convolution variants, avoiding additional dedicated arithmetic units. The TCE consists of control logic, a shared multiply–accumulate array, an adder tree, and a multifunction activation unit. The datapath undergoes runtime reconfiguration when the workload transitions to the multibranch gating stage. During this phase, concurrent access by the main branch, channel gate, and spatial gate exceeds the port limits of standard dual-port BRAMs, leading to structural hazards. We address this by implementing a multibank buffer that supports one write stream and three read streams per pipeline stage. To further reduce pipeline depth, the output stage incorporates fused logic that performs batch normalization, quantization, and activation within a single datapath.

The division-free encoding module connects the synchronous numeric and asynchronous event domains by eliminating the floating-point division used in conventional TTFS encoders. By constraining the scaling factor $S_p$ to a power of two ($S_p=2^k$) during quantization-aware training, a barrel shifter performs the division through a right shift. As illustrated in Fig.~\ref{fig:hardoverview}c, the encoder hardware employs bit-shift normalization and event detection logic to generate AER spike packets directly from INT8 features. The resulting combinational datapath computes spike latencies within a single clock cycle.

To address frequent weight access and sparse neural activity in SNN inference, the compute engine adopts a memory-centric array architecture. Co-locating weight storage with computation alleviates the bandwidth limitations of the conventional von Neumann architecture and reduces data movement overhead. The compute fabric is partitioned into 16 independent SNN arrays, each containing four event-driven processing elements. The microarchitecture of each PE, shown in Fig.~\ref{fig:hardoverview}d, performs time-modulated membrane potential integration and is controlled by a local asynchronous handshake mechanism. The AER interface manages broadcasting and recirculation, connecting sparse event streams to parallel arrays while providing a bidirectional datapath for time-multiplexed layer execution.

To evaluate hardware deployment, we implemented the INT8 model on a Xilinx Zynq-7020 SoC. The accelerator resides in the programmable logic and operates at 100~MHz with a 1.0~V core supply voltage. Resource utilization confirms the compact hardware footprint of the design (Table~S3). Look-up table and flip-flop utilization remain low at 20.34\% and 7.67\%, respectively, leaving margin for future architectural scaling. The implementation achieves a latency of 0.018~ms for a single EEG sliding-window inference. The accelerator consumes 40~mW of dynamic power and 108~mW of static power, resulting in a total power consumption of 148~mW (Table~S4). This yields an energy consumption of 2.66~$\mu$J per inference.


\section{DISCUSSION}
Overall, the results suggest that the strength of the proposed framework lies in cross-layer co-design rather than in any single architectural component \cite{ref71}. Across five EEG emotion-recognition benchmarks, the framework maintains strong performance under different emotional granularities and subject conditions while remaining robust to moderate noise perturbations and 8-bit quantization. These observations indicate that sparse temporal coding can preserve affect-relevant information and improve deployability, rather than merely optimizing a single benchmark setting.

This behavior can be explained by the complementary roles of the main algorithmic components. DSGM does more than reduce parameters: by jointly reweighting spatial locations and recording channels, it concentrates discriminative responses and suppresses irrelevant activations. DF-TTFS then converts dense features into sparse spike times using power-of-two scaling, eliminating costly division while preserving the ordering of salient responses. ATSNN further stabilizes temporal learning by adapting the integration window to the earliest valid spike, thereby mitigating the skewed presynaptic offsets that often lead to unstable gradients in TTFS-based training \cite{ref70}. Together, these mechanisms enable sparse event generation without a substantial loss of representational fidelity.

The efficiency results also show that the benefits of sparsity extend beyond the algorithmic level. On Raspberry Pi 5, sparse firing reduces effective operations and memory demand, allowing the quantized network to operate within a compact storage and runtime footprint. On hardware, the GALS accelerator converts the same sparsity into measurable system-level gains by activating communication, memory access, and computation only when valid events are present. This event-proportional execution avoids much of the synchronization overhead associated with clock-driven designs and demonstrates that the theoretical efficiency of SNNs can be translated into practical energy savings.

Several limitations remain. The present study is based on public benchmarks, controlled noise perturbations, and an FPGA prototype, whereas long-term wearable deployment will face additional variability arising from electrode impedance drift, motion artifacts, and subject-specific state changes. Future work should therefore investigate adaptive event-generation thresholds, online calibration strategies, and tighter integration between front-end acquisition and event-driven inference. In addition, migrating the current programmable-logic prototype to an ASIC implementation would help clarify the achievable area and power envelope for continuous wearable affective monitoring.

\section{CONCLUSION}
This work presents an end-to-end neuromorphic co-design for low-power wearable EEG emotion recognition. By combining sparse temporal encoding, spiking inference, and asynchronous hardware execution, the proposed framework preserves strong recognition performance across five benchmark datasets while remaining robust to moderate noise and 8-bit quantization. These results show that affective EEG can be processed as sparse events without sacrificing practical recognition capability, providing a deployment-oriented alternative to synchronous dense computation.

Efficiency gains are observed at both the model and hardware levels. The INT8 network requires only 12.50~KB of parameter storage and 13.68~KB of runtime memory and consumes 261.21~$\mu$J per inference on Raspberry Pi~5. The FPGA implementation of the GALS accelerator further achieves 0.018~ms latency, 40~mW dynamic power, and 2.66~$\mu$J per inference. Together, these findings demonstrate that algorithmic sparsity can be converted into measurable system-level efficiency gains, highlighting asynchronous neuromorphic computing as a practical direction for wearable affective interfaces.

\section{METHODS}
Detailed methods are given in the online Supplementary data.

\section{DATA AND CODE AVAILABILITY}
The source code is available at \url{https://github.com/BhealthE6315/DA-SNN}.

\section{FUNDING}
This work was supported by the National Natural Science Foundation of China (No. 62227807).

\bibliographystyle{nsr}
\bibliography{nsr_sample}
\end{document}